\section{Conclusion}

Purpose-selective execution asks for more than hiding data or deleting it: a
confidential attribute should remain useful where approved and irrelevant where
prohibited. This paper defines a bounded way to test that requirement across
data, workloads, capabilities, releases, privacy, lifecycle, and evidence.

The diagnostic evaluation demonstrates that the PSBE-Runtime pipeline executes
correctly against three commercial models via OpenRouter. The admission gate
rejects models with unstable route assignments, the position diagnostic
identifies layout-sensitive release behavior, and the R2 diagnostic
validates end-to-end execution across datasets and conditions. Release
denials are concentrated in one model on one dataset.

A reduced live lane (330 invocations of one admitted model) is reported
strictly as a live execution and release-policy diagnostic over one
commercial model lane. An audit of its transmitted-field evidence found that
only a pseudonymized record identifier was transmitted, so the stream carries
no purpose-selectivity contrast; the confirmatory experiment is suspended
until that defect is corrected and re-validated.

Attack, differential-privacy, evidence-verification, availability, and
overhead results are intentionally withheld: prior test-double values have
been withdrawn and will be reported only from live attacks, real DP-SGD
training, real evidence verification, and measured wall-clock executions.

Three diagnostic claims are supported by the live diagnostic evidence. The
full confirmatory claim family (H1--H10) requires the live non-TEE
purpose-selective benchmark with statistical inference.
